\begin{textsample}{POS Dim 4 – rn – Score 11.00 – rn/rn\_em\_15.txt}  \label{ex:f4_pos_053}
3.6 \textbf{POPULAÇÕES} EM SITUAÇÃO DE ITINERÂNCIA A educação das \textbf{populações} em situação de itinerância, embora não conste como modalidade de ensino, seu atendimento educacional é legítimo pelo processo histórico de luta e resistência dessas \textbf{populações}, seus valores étnico-raciais, políticos e culturais.
Desse modo, neste Referencial Curricular, propomos o atendimento educacional diferenciado a esses grupos.
Nessa perspectiva, os espaços, o currículo e as vivências pedagógicas do atendimento educacional dessas \textbf{populações} se fundamentam no reconhecimento e na valorização da diversidade cultural, na sua relação com a terra e o trabalho, modo de organização coletiva e saberes.
A Coordenação-Geral de Educação para as Relações Étnico-Raciais, conforme Resolução CNE/CEB nº 3, de 16 de maio de 2012, com base no Parecer CNE/CEB nº 14/2011, definiu e situou as \textbf{populações} em situação de itinerância, igualmente, determinou que os sistemas de ensino deverão adequar -se às particularidades dos estudantes, garantindo, assim, seus direitos socioeducacionais.
Desse modo, a Rede Estadual de Ensino do RN atende à legislação no que diz respeito à oferta de educação às \textbf{populações} em situação de itinerância, conforme disposto no parágrafo único do art. 1º da Resolução nº 3, de 16 de maio de 2012 : > Art. 1º As crianças, adolescentes e jovens em situação de itinerância deverão ter garantido o direito à matrícula em escola pública, gratuita, com qualidade social e que garanta a liberdade de consciência e de crença. > > Parágrafo único.
São considerados crianças, adolescentes e jovens em situação de itinerância aquelas pertencentes a grupos sociais que vivem em tal condição por motivos culturais, políticos, econômicos, de saúde, tais como \textbf{ciganos}, \textbf{indígenas}, povos nômades, trabalhadores itinerantes, acampados, circenses, \textbf{artistas} e/ou trabalhadores de parques de diversão, de teatro mambembe, dentre outros.
Nesse intento, o ensino médio do RN se organiza a partir das orientações contidas neste Referencial, direcionando o planejamento de ações pedagógicas nas escolas, voltadas ao atendimento às \textbf{populações} consideradas em situação de itinerância que se estabeleçam no \textbf{território} potiguar.
A condição humana de muitos destes grupos é de poucas oportunidades de assistência educacional, em decorrência da vivência em contextos diversos, ocasionada pelos constantes deslocamentos.
Entre as \textbf{populações} em situação de itinerância, destacam -se os \textbf{ciganos}, por constituírem grupos representativos no nosso \textbf{território}, sendo imperativo pensarmos e executarmos ações educativas voltadas a esses povos.
Segundo dados do Instituto Brasileiro de Geografia e Estatística ( IBGE ), em 2011, a \textbf{população} \textbf{cigana} total constava de cerca de 800.000 \textbf{ciganos} no Brasil, com acampamentos distribuídos em 21 estados da federação.
Para Silva ( 2016, p.60 ), “ sempre houve uma circularidade de grupos de \textbf{ciganos} entre os estados do \textbf{nordeste} brasileiro e, em especial, entre os estados da Paraíba, Rio Grande do Norte e Ceará ”.
Sendo que no \textbf{território} potiguar registra -se a existência de grupos de \textbf{ciganos} denominados de Calons, Matchuanos, Roraronês e Kalderash.
Em coerência com a declaração de Silva ( 2016 ), os dados da Pesquisa de Informações Básicas Municipais ( MUNIC/IBGE – 2011 ) e declaração de \textbf{ciganos} participantes do GT-MEC-CIGANOS, indicam que há no Rio Grande do Norte, a presença de acampamentos \textbf{ciganos} nos seguintes municípios : Apodi, Bom Jesus, Carnaúba dos Dantas, Equador, Florânia, Luís Gomes, Macaíba, Macau, Mossoró, Parnamirim, Pau dos Ferros, Presidente Juscelino, São Paulo do Potengi, São Vicente, Santa Cruz e Tangará.
Compreendemos que a itinerância é condição que impacta a matrícula escolar e a permanência dos estudantes, pertencentes a esses grupos sociais, no percurso do ensino médio.
Nesse sentido, é necessário que a Rede Estadual de Ensino normatize, assegure a matrícula e ofereça, gradativamente, as condições essenciais para que essa \textbf{população} de itinerantes frequentem as escolas potiguares, de modo a minimizar a descontinuidade na aprendizagem, garantindo acesso à educação enquanto direito social, com liberdade de crença, sem qualquer forma de preconceito e/ou qualquer forma de discriminação.
Portanto, é fundamental que as unidades escolares, que acolhem os estudantes dessa \textbf{população}, contemplem em seus Projetos Pedagógicos ações no que concerne o atendimento a esses jovens, que propiciem a formação integral, que discutam sobre temas importantes para estes grupos, como : direito à educação, cidadania, direitos humanos, de acesso à terra e à moradia, dentre outros.
Por fim, ressalta -se a importância de as escolas promoverem debates relacionados à cultura, modo de vida e compreensão das questões étnicos-raciais que perpassam vários grupos com as quais os estudantes têm acesso ou estabelecem relações.

% matched lemmas: artista, cigano, indígena, nordeste, população, território
\end{textsample}
